\documentclass[11pt]{article}

\usepackage[margin=1in]{geometry}
\usepackage{amsmath, amssymb, amsthm}
\usepackage{enumitem}

\newcommand{\Z}{\mathbb{Z}}
\newcommand{\R}{\mathbb{R}}

\pagestyle{empty}

\begin{document}

\begin{center}
  {\Large\bfseries Weekly Assignment: Counting and Probability (Epp Ch.\ 9)}
\end{center}

\bigskip

\noindent
Complete each of the following problems.  Show all work and justify your answers.

\bigskip

%% ─────────────────────────────────────────────
\noindent\textbf{Problem 1} (Epp 9.1 \#4). \;
Use the sample space given in Example~9.1.1 (a standard 52-card deck).
Write the following event as a set and compute its probability:

\medskip
\emph{The event that the chosen card is black and has an even number on it.}

\bigskip

%% ─────────────────────────────────────────────
\noindent\textbf{Problem 2} (Epp 9.1 \#10). \;
Use the sample space given in Example~9.1.2 (rolling a blue die and a gray die
together).  Write the following event as a set and compute its probability:

\medskip
\emph{The event that the sum of the numbers showing face up is at least~9.}

\bigskip

%% ─────────────────────────────────────────────
\noindent\textbf{Problem 3} (Epp 9.1 \#19). \;
An urn contains two blue balls (denoted $B_1$ and $B_2$) and three white balls
(denoted $W_1$, $W_2$, and $W_3$).  One ball is drawn, its color is recorded,
and it is replaced in the urn.  Then another ball is drawn and its color is recorded.
Let $B_1 W_2$ denote the outcome that the first ball drawn is $B_1$ and the
second ball drawn is $W_2$.  Because the first ball is replaced before the second
ball is drawn, the outcomes of the experiment are equally likely.
\begin{enumerate}[label=\alph*.]
  \item List all 25 possible outcomes of the experiment.
  \item Consider the event that the first ball drawn is blue.
        List all outcomes in the event.  What is the probability of the event?
  \item Consider the event that only white balls are drawn.
        List all outcomes in the event.  What is the probability of the event?
\end{enumerate}

\bigskip

%% ─────────────────────────────────────────────
\noindent\textbf{Problem 4} (Epp 9.2 \#10). \;
Suppose there are three routes from North Point to Boulder Creek, two routes
from Boulder Creek to Beaver Dam, two routes from Beaver Dam to Star Lake,
and four routes directly from Boulder Creek to Star Lake.
\begin{enumerate}[label=\alph*.]
  \item How many routes from North Point to Star Lake pass through Beaver Dam?
  \item How many routes from North Point to Star Lake bypass Beaver Dam?
\end{enumerate}

\bigskip

%% ─────────────────────────────────────────────
\noindent\textbf{Problem 5} (Epp 9.2 \#26). \;
Determine how many times the innermost loop will be iterated when the
following algorithm segment is implemented and run.
(Assume that $m$, $n$, and $p$ are all positive integers.)

\medskip
\begin{tabbing}
\quad \= \quad \= \quad \= \kill
\textbf{for} $i := 1$ \textbf{to} $m$ \\
\> \textbf{for} $j := 1$ \textbf{to} $n$ \\
\> \> \textbf{for} $k := 1$ \textbf{to} $p$ \\
\> \> \> [Statements in body of inner loop.] \\
\> \> \> [None contain branching statements that lead outside the loop.] \\
\> \> \textbf{next} $k$ \\
\> \textbf{next} $j$ \\
\textbf{next} $i$
\end{tabbing}

\bigskip

%% ─────────────────────────────────────────────
\noindent\textbf{Problem 6} (Epp 9.3 \#8). \;
In a certain country, license plates consist of zero or one digit followed by
four or five uppercase letters from the Roman alphabet.
\begin{enumerate}[label=\alph*.]
  \item How many different license plates can the country produce?
  \item How many license plates have no repeated letter?
  \item How many license plates have at least one repeated letter?
  \item What is the probability that a license plate has a repeated letter?
\end{enumerate}

\bigskip

%% ─────────────────────────────────────────────
\noindent\textbf{Problem 7} (Epp 9.4 \#6).
\begin{enumerate}[label=\alph*.]
  \item Given any set of seven integers, must there be two that have the same
        remainder when divided by~6?  Why?
  \item Given any set of seven integers, must there be two that have the same
        remainder when divided by~8?  Why?
\end{enumerate}

\bigskip

%% ─────────────────────────────────────────────
\noindent\textbf{Problem 8}. \;
Consider strings of length $n$ over the set $\{1, 2, 3, 4\}$.  For example,
one element of the set of strings of length~5 would be $12432$.

\medskip
What is the probability that a string of length $n$ has at least two
\emph{adjacent} characters that are the same (as in $14\mathbf{33}1$)?

\medskip\noindent
\emph{Hint}: It may be easier to first count the strings with \textbf{no} two
adjacent characters equal, and then use the complement.

\end{document}
